\documentclass[../../main.tex]{subfiles}
\begin{document}
\section{Riemann 流形}\label{section:ch13-riemannian-manifolds}

向量空间上最重要的对称张量例子就是内积。任何内积都使我们能够定义向量的长度以及它们之间的夹角，从而进行 Euclid 几何。

把这些想法移植到流形上，我们便得到张量在微分几何中最重要的应用之一。设 $M$ 是一个带边或不带边的光滑流形。$M$ 上的一个 \emph{Riemann 度量} 是 $M$ 上一个光滑对称协变 $2$-张量场，它在每一点处都是正定的。一个 \emph{Riemann 流形} 是一对 $(M,g)$，其中 $M$ 是一个光滑流形，$g$ 是 $M$ 上的一个 Riemann 度量。当 $M$ 被理解为已被赋予某个特定的 Riemann 度量时，人们有时径直说“$M$ 是一个 Riemann 流形”。带边 Riemann 流形可类似地定义。

注意，Riemann 度量与度量空间意义下的度量并不是一回事，尽管这两个概念是相关的，我们将在下面看到。由于这种歧义，我们通常用“距离函数”一词来指度量空间意义下的度量，而把“度量”一词留给 Riemann 度量。无论如何，所讨论的是哪一种度量，总应能从上下文中看清楚。

若 $g$ 是 $M$ 上的一个 Riemann 度量，则对每个 $p\in M$，$2$-张量 $g_p$ 是 $T_pM$ 上的一个内积。正因如此，我们常采用记号 $\langle v,w\rangle_g$ 来表示实数 $g_p(v,w)$，其中 $v,w\in T_pM$。

在任意光滑局部坐标 $(x^i)$ 下，Riemann 度量可写成
\begin{align*}
g=g_{ij}\,\mathrm{d}x^i\otimes \mathrm{d}x^j,
\end{align*}
其中 $(g_{ij})$ 是一个由光滑函数组成的对称正定矩阵。由 $g$ 的对称性，我们还可以用对称乘积把它写成如下形式：
\begin{align*}
g&=g_{ij}\,\mathrm{d}x^i\otimes \mathrm{d}x^j\\
&=\tfrac12\bigl(g_{ij}\,\mathrm{d}x^i\otimes \mathrm{d}x^j+g_{ji}\,\mathrm{d}x^i\otimes \mathrm{d}x^j\bigr)
\quad(\text{因为 }g_{ij}=g_{ji})\\
&=\tfrac12\bigl(g_{ij}\,\mathrm{d}x^i\otimes \mathrm{d}x^j+g_{ij}\,\mathrm{d}x^j\otimes \mathrm{d}x^i\bigr)
\quad(\text{在第二项中交换 }i\leftrightarrow j)\\
&=g_{ij}\,\mathrm{d}x^i\mathrm{d}x^j
\quad(\text{由\cref{proposition:ch12-symmetric-product}(b)}).
\end{align*}

\begin{example}[Euclid 度量]\label{example:ch13-euclidean-metric}
Riemann 度量最简单的例子是 $\mathbb{R}^n$ 上的 Euclid 度量 $\bar g$，它在标准坐标下由
\begin{align*}
\bar g=\delta_{ij}\,\mathrm{d}x^i\mathrm{d}x^j
\end{align*}
给出，其中 $\delta_{ij}$ 是 Kronecker 符号。人们通常把一个张量 $\alpha$ 与自身的对称乘积缩写为 $\alpha^2$，于是 Euclid 度量也可写成
\begin{align*}
\bar g=(\mathrm{d}x^1)^2+\cdots+(\mathrm{d}x^n)^2.
\end{align*}
把它作用于向量 $v,w\in T_p\mathbb{R}^n$，便得到
\begin{align*}
\bar g_p(v,w)=\delta_{ij}v^iw^j=\sum_{i=1}^n v^iw^j=v\cdot w.
\end{align*}
换言之，$\bar g$ 是这样一个 $2$-张量场：它在每一点处的值就是 Euclid 点积。（你也许还记得，我们在第~1 章中曾提醒过，涉及 Euclid 点积的表达式很容易违反我们的指标约定，因而需要显式的求和号。通过像 $\delta_{ij}v^iw^j$ 这样显式写出度量系数 $\delta_{ij}$，通常就可以避免这一点。）
\end{example}

\begin{example}[乘积度量]\label{example:ch13-product-metrics}
若 $(M,g)$ 与 $(\widetilde M,\widetilde g)$ 是 Riemann 流形，则我们可以在乘积流形 $M\times\widetilde M$ 上定义一个 Riemann 度量 $\hat g=g\oplus\widetilde g$，称为\emph{乘积度量}，定义如下：
\begin{align}
\hat g\bigl((v,\widetilde v),(w,\widetilde w)\bigr)=g(v,w)+\widetilde g(\widetilde v,\widetilde w)
\end{align}
对任意 $(v,\widetilde v),(w,\widetilde w)\in T_pM\oplus T_q\widetilde M\cong T_{(p,q)}(M\times\widetilde M)$ 成立。给定 $M$ 的任意局部坐标 $(x^1,\dots,x^n)$ 和 $\widetilde M$ 的任意局部坐标 $(y^1,\dots,y^m)$，我们便得到 $M\times\widetilde M$ 的局部坐标 $(x^1,\dots,x^n,y^1,\dots,y^m)$，你可以验证，乘积度量局部地由分块对角矩阵
\begin{align*}
(\hat g_{ij})=
\begin{pmatrix}
g_{ij}&0\\
0&\widetilde g_{ij}
\end{pmatrix}
\end{align*}
表示。例如，容易验证 $\mathbb{R}^{n+m}$ 上的 Euclid 度量与由 $\mathbb{R}^n$ 和 $\mathbb{R}^m$ 上的 Euclid 度量所确定的乘积度量相同。
\end{example}

Riemann 度量一个令人愉快的特点是它们大量存在。（另一种处理方式见习题\ref{exercise:ch13-18}。）

\begin{proposition}[Riemann 度量的存在性]\label{proposition:ch13-existence}
每个带边或不带边的光滑流形都容许一个 Riemann 度量。
\end{proposition}

\begin{proof}
设 $M$ 是一个带边或不带边的光滑流形，选取 $M$ 的一个由光滑坐标图 $(U_\alpha,\varphi_\alpha)$ 组成的覆盖。在每个坐标域上，有一个 Riemann 度量 $g_\alpha=\varphi_\alpha^*\bar g$，其坐标表达式为 $\delta_{ij}\,\mathrm{d}x^i\mathrm{d}x^j$。令 $\{\varphi_\alpha\}$ 是从属于覆盖 $\{U_\alpha\}$ 的一个光滑单位分解，并定义
\begin{align*}
g=\sum_\alpha \varphi_\alpha g_\alpha,
\end{align*}
其中每一项在 $\operatorname{supp}\varphi_\alpha$ 之外被理解为零。由局部有限性，在每一点的邻域内只有有限多个非零项，因此这个表达式定义了一个光滑张量场。它显然是对称的，所以只需验证正定性。若 $v\in T_pM$ 是任意非零向量，则
\begin{align*}
g_p(v,v)=\sum_\alpha \varphi_\alpha(p)g_{\alpha,p}(v,v).
\end{align*}
这个和是非负的，因为每一项都是非负的。函数 $\varphi_\alpha$ 中至少有一个在 $p$ 处严格为正（因为它们之和为 $1$）。又因为 $g_{\alpha,p}(v,v)>0$，便得 $g_p(v,v)>0$。
\end{proof}

需要着重指出：对于给定的流形，构造一个度量 $g$ 时有极大的选择余地，因此它没有任何典范性。特别地，同一流形上不同的度量可以具有大不相同的几何性质。例如，习题\ref{exercise:ch13-20}描述了 $\mathbb{R}^2$ 上四个行为迥然不同的度量。

下面只是可在带边或不带边的 Riemann 流形 $(M,g)$ 上定义的若干几何构造中的几个。

\begin{itemize}
\item 切向量 $v\in T_pM$ 的\emph{长度}或\emph{范数}定义为
\begin{align*}
|v|_g=\langle v,v\rangle_g^{1/2}=g_p(v,v)^{1/2}.
\end{align*}
\item 两个非零切向量 $v,w\in T_pM$ 之间的\emph{夹角}是满足下式的唯一 $\theta\in[0,\pi]$：
\begin{align*}
\cos\theta=\frac{\langle v,w\rangle_g}{|v|_g|w|_g}.
\end{align*}
\item 若 $\langle v,w\rangle_g=0$，则称切向量 $v,w\in T_pM$ 是\emph{正交的}。这意味着其中一个或两个向量为零，或者它们之间的夹角为 $\pi/2$。
\end{itemize}

研究 Riemann 流形时一个极为有用的工具是标准正交标架。设 $(M,g)$ 是一个 $n$ 维带边或不带边的 Riemann 流形。正如我们在第~8 章中对 $\mathbb{R}^n$ 所做的那样，若 $M$ 在开子集 $U\subseteq M$ 上的一个局部标架 $(E^1,\dots,E^n)$ 在每一点 $p\in U$ 处使向量 $(E^1|_p,\dots,E^n|_p)$ 构成 $T_pM$ 的一组标准正交基，等价地说若 $\langle E^i,E^j\rangle_g=\delta^i_j$，则称它为一个\emph{标准正交标架}。

\begin{example}\label{example:ch13-coord-frame-orthonormal}
坐标标架 $(\partial/\partial x^i)$ 是带 Euclid 度量的 $\mathbb{R}^n$ 上的一个整体标准正交标架。
\end{example}

\begin{example}\label{example:ch13-polar-frame}
在例\ref{example:ch08-polar-frame}中定义于 $\mathbb{R}^2\setminus\{0\}$ 上的标架 $(E^1,E^2)$ 是 $\mathbb{R}^2$ 的一个局部标准正交标架。正如我们在例\ref{example:ch09-9-45}中所观察到的，它在任何坐标下都不是坐标标架。
\end{example}

下一个命题的证明与引理\ref{lemma:ch08-gram-schmidt}的证明完全相同，只需把 Euclid 点积换成内积 $\langle\cdot,\cdot\rangle_g$。

\begin{proposition}\label{proposition:ch13-gram-schmidt}
设 $(M,g)$ 是一个带边或不带边的 Riemann 流形，$(X^j)$ 是 $M$ 在开子集 $U\subseteq M$ 上的一个光滑局部标架。则存在 $U$ 上的一个光滑标准正交标架 $(E^j)$，使得对每个 $j=1,\dots,n$ 和每个 $p\in U$，
\begin{align*}
\operatorname{span}(E^1|_p,\dots,E^j|_p)=\operatorname{span}(X^1|_p,\dots,X^j|_p).
\end{align*}
\end{proposition}

\begin{exercise}\label{exercise:ch13-07}
证明上一命题。
\end{exercise}

\begin{corollary}[局部标准正交标架的存在性]\label{corollary:ch13-local-orthonormal}
设 $(M,g)$ 是一个带边或不带边的 Riemann 流形。对每个 $p\in M$，在 $p$ 的某个邻域上存在一个光滑标准正交标架。
\end{corollary}

\begin{proof}
从一个光滑坐标标架出发，应用命题\ref{proposition:ch13-gram-schmidt}即可。
\end{proof}

注意，推论\ref{corollary:ch13-local-orthonormal}并未表明在 $p$ 的邻域上存在使坐标标架标准正交的光滑坐标。事实上，正如我们将在下面看到的，情形很少如此。

\subsection{拉回度量}

设 $M,N$ 是带边或不带边的光滑流形，$g$ 是 $N$ 上的一个 Riemann 度量，$F:M\to N$ 光滑。拉回 $F^*g$ 是 $M$ 上的一个光滑 $2$-张量场。若它是正定的，则它是 $M$ 上的一个 Riemann 度量，称为由 $F$ 确定的\emph{拉回度量}。下一个命题说明何时如此。

\begin{proposition}[拉回度量判别法]\label{proposition:ch13-pullback-criterion}
设 $F:M\to N$ 是一个光滑映射，$g$ 是 $N$ 上的一个 Riemann 度量。则 $F^*g$ 是 $M$ 上的 Riemann 度量当且仅当 $F$ 是一个光滑浸入。
\end{proposition}

\begin{exercise}\label{exercise:ch13-10}
证明上一命题。
\end{exercise}

若已知一个浸入的坐标表示，则用通常计算拉回的算法很容易算出拉回度量。

\begin{example}\label{example:ch13-helicoid}
考虑由 $F(u,v)=(u\cos v,u\sin v,v)$ 给出的光滑映射 $F:\mathbb{R}^2\to\mathbb{R}^3$。它是一个正常的单射光滑浸入，因而由命题\ref{proposition:ch04-embedding-criterion}是一个嵌入。它的像就是一张称为\emph{螺旋面}的曲面；它看起来像一条无限宽的螺旋形斜坡（图\ref{figure:ch13-helicoid}）。拉回度量 $F^*\bar g$ 可以通过把 $F$ 的坐标函数代入 $\bar g$ 公式中的 $x,y,z$ 而算出：
\begin{align*}
F^*\bar g&=\mathrm{d}(u\cos v)^2+\mathrm{d}(u\sin v)^2+\mathrm{d}(v)^2\\
&=(\cos v\,\mathrm{d}u-u\sin v\,\mathrm{d}v)^2+(\sin v\,\mathrm{d}u+u\cos v\,\mathrm{d}v)^2+\mathrm{d}v^2\\
&=\cos^2 v\,\mathrm{d}u^2-2u\sin v\cos v\,\mathrm{d}u\,\mathrm{d}v+u^2\sin^2 v\,\mathrm{d}v^2\\
&\quad+\sin^2 v\,\mathrm{d}u^2+2u\sin v\cos v\,\mathrm{d}u\,\mathrm{d}v+u^2\cos^2 v\,\mathrm{d}v^2+\mathrm{d}v^2\\
&=\mathrm{d}u^2+(u^2+1)\,\mathrm{d}v^2.
\end{align*}
（按惯例，当 $u$ 是一个实值函数时，记号 $\mathrm{d}u^2$ 指对称乘积 $\mathrm{d}u\,\mathrm{d}u$，而不是 $\mathrm{d}(u^2)$。）
\end{example}

\begin{figure}[htbp]
\centering
\includegraphics[width=0.5\textwidth]{fig13-1-helicoid.png}
\caption{一张螺旋面}
\label{figure:ch13-helicoid}
\end{figure}

为在坐标变换下变换 Riemann 度量，我们使用与处理余向量场时相同的技巧：把坐标变换视为定义域与陪域采用不同坐标的恒等映射，并用推论\ref{corollary:ch12-pullback-tensor}的公式来计算拉回。如前所述，在实践中这只不过是把一组坐标用另一组坐标表示的公式代入而已。

\begin{example}\label{example:ch13-polar-metric}
为说明这一点，我们在极坐标下计算 $\mathbb{R}^2$ 上 Euclid 度量 $\bar g=\mathrm{d}x^2+\mathrm{d}y^2$ 的坐标表达式。代入 $x=r\cos\theta$、$y=r\sin\theta$ 并展开，得到
\begin{align*}
\bar g&=\mathrm{d}x^2+\mathrm{d}y^2=\mathrm{d}(r\cos\theta)^2+\mathrm{d}(r\sin\theta)^2\\
&=(\cos\theta\,\mathrm{d}r-r\sin\theta\,\mathrm{d}\theta)^2+(\sin\theta\,\mathrm{d}r+r\cos\theta\,\mathrm{d}\theta)^2\\
&=(\cos^2\theta+\sin^2\theta)\,\mathrm{d}r^2+(r^2\sin^2\theta+r^2\cos^2\theta)\,\mathrm{d}\theta^2\\
&\quad+(-2r\cos\theta\sin\theta+2r\sin\theta\cos\theta)\,\mathrm{d}r\,\mathrm{d}\theta\\
&=\mathrm{d}r^2+r^2\,\mathrm{d}\theta^2.
\end{align*}
\end{example}

若 $(M,g)$ 与 $(\widetilde M,\widetilde g)$ 都是 Riemann 流形，则一个光滑映射 $F:M\to\widetilde M$ 若为微分同胚且满足 $F^*\widetilde g=g$，就称为一个（Riemann）\emph{等距}。更一般地，若每一点 $p\in M$ 都有一个邻域 $U$ 使 $F|_U$ 是 $U$ 到 $\widetilde M$ 的某个开子集上的等距，则称 $F$ 为一个\emph{局部等距}；等价地说，若 $F$ 是满足 $F^*\widetilde g=g$ 的局部微分同胚。

若在 $(M,g)$ 与 $(\widetilde M,\widetilde g)$ 之间存在一个 Riemann 等距，我们就说它们作为 Riemann 流形是\emph{等距的}。若 $M$ 的每一点都有一个邻域与 $(\widetilde M,\widetilde g)$ 的某个开子集等距，则称 $(M,g)$\emph{局部等距于} $(\widetilde M,\widetilde g)$。研究 Riemann 流形在（局部或整体）等距下不变的性质，称为 Riemann 几何。

其中一个这样的性质是平坦性。若 $(M,g)$ 局部等距于 $(\mathbb{R}^n,\bar g)$，则称这个 $n$ 维 Riemann 流形 $(M,g)$ 为\emph{平坦 Riemann 流形}，而 $g$ 为\emph{平坦度量}。

\begin{exercise}\label{exercise:ch13-13}
设 $(M,g)$ 与 $(\widetilde M,\widetilde g)$ 是等距的 Riemann 流形。证明 $g$ 平坦当且仅当 $\widetilde g$ 平坦。
\end{exercise}

下一个定理是判定一个 Riemann 度量是否平坦的关键。

\begin{theorem}\label{theorem:ch13-flatness}
对于一个 Riemann 流形 $(M,g)$，下列各条等价：
\begin{enumerate}[(a)]
\item $g$ 是平坦的。
\item $M$ 的每一点都含于某个光滑坐标图的定义域内，在其中 $g$ 具有坐标表示 $g=\delta_{ij}\,\mathrm{d}x^i\mathrm{d}x^j$。
\item $M$ 的每一点都含于某个光滑坐标图的定义域内，在其中坐标标架是标准正交的。
\item $M$ 的每一点都含于某个交换标准正交标架的定义域内。
\end{enumerate}
\end{theorem}

\begin{proof}
蕴含关系 (a)$\Rightarrow$(b)$\Rightarrow$(c)$\Rightarrow$(d) 是定义的直接推论，留给读者。余下的蕴含关系 (d)$\Rightarrow$(a) 来自交换标架的典范形式定理：若 $(E^i)$ 是 $g$ 在开子集 $U\subseteq M$ 上的一个交换标准正交标架，则定理\ref{theorem:ch09-commuting-frame}表明每一个 $p\in V$ 都含于某个光滑坐标图 $(U,\varphi)$ 的定义域内，在其中坐标标架等于 $(E^i)$。这意味着 $\varphi_*E^i=\partial/\partial x^i$，于是微分同胚 $\varphi:U\to\varphi(U)$ 满足
\begin{align*}
\varphi^*\bar g(E^i,E^j)=\bar g(\varphi_*E^i,\varphi_*E^j)
=\bar g\!\left(\frac{\partial}{\partial x^i},\frac{\partial}{\partial x^j}\right)
=\delta^i_j=g(E^i,E^j).
\end{align*}
双线性性随即给出 $\varphi^*\bar g=g$，故 $\varphi$ 是 $(U,g|_U)$ 与带 Euclid 度量的 $\varphi(U)$ 之间的一个等距。这表明 $g$ 是平坦的。
\end{proof}

\begin{exercise}\label{exercise:ch13-15}
通过证明 (a)$\Rightarrow$(b)$\Rightarrow$(c)$\Rightarrow$(d) 来补全上一证明。
\end{exercise}

从定义本身完全看不出存在非平坦的 Riemann 度量。事实上，在一维情形，每个度量都是平坦的，正如习题\ref{exercise:ch13-06}所示。在本章后面，我们将用定理\ref{theorem:ch13-flatness}证明 $\mathbb{R}^3$ 中大多数旋转曲面（包括 $S^2$）都不是平坦的。

\subsection{Riemann 子流形}

拉回度量对子流形尤为重要。若 $(M,g)$ 是一个带边或不带边的 Riemann 流形，则每个子流形 $S\subseteq M$（浸入或嵌入，带边或不带边）都自动继承一个拉回度量 $\iota^*g$，其中 $\iota:S\hookrightarrow M$ 是包含映射。在这一背景下，拉回度量也称为 $S$ 上的\emph{诱导度量}。按定义，这意味着对 $v,w\in T_pS$，
\begin{align*}
(\iota^*g)(v,w)=g(\mathrm{d}\iota_p(v),\mathrm{d}\iota_p(w))=g(v,w),
\end{align*}
因为 $\mathrm{d}\iota_p:T_pS\to T_pM$ 就是我们通常把 $T_pS$ 等同于 $T_pM$ 的子空间的那个映射。因此 $\iota^*g$ 只不过是 $g$ 在与 $S$ 相切的向量对上的限制。带上这个度量后，$S$ 称为 $M$ 的一个\emph{Riemann 子流形}（带边或不带边）。

\begin{example}\label{example:ch13-round-metric}
由通常的包含映射 $\iota:S^n\hookrightarrow\mathbb{R}^{n+1}$ 在 $S^n$ 上诱导的度量 $\bar g=\iota^*\bar g$，称为球面上的\emph{圆度量}（或\emph{标准度量}）。
\end{example}

若 $(M,g)$ 是一个 Riemann 流形，$\iota:S\hookrightarrow M$ 是一个 Riemann 子流形，则计算诱导度量 $\iota^*g$ 通常最容易借助局部参数化来进行；回忆第~5 章，这是指一个从开子集 $U\subseteq\mathbb{R}^k$ 到 $M$ 的单射浸入 $X$，其像是 $S$ 的一个开子集，而其逆是 $S$ 的一个光滑坐标映射。由于 $\iota\circ X=X$，$\iota^*g$ 的坐标表示为 $X^*(\iota^*g)=X^*g$。下面两个例子说明这一过程。

\begin{example}[图坐标中的诱导度量]\label{example:ch13-graph-coords}
设 $U\subseteq\mathbb{R}^n$ 是一个开子集，$S\subseteq\mathbb{R}^{n+1}$ 是一个光滑函数 $f:U\to\mathbb{R}$ 的图像。由 $X(u^1,\dots,u^n)=(u^1,\dots,u^n,f(u))$ 给出的映射 $X:U\to\mathbb{R}^{n+1}$ 是 $S$ 的一个光滑整体参数化，$S$ 上的诱导度量在图坐标下由
\begin{align*}
X^*\bar g&=X^*\bigl((\mathrm{d}x^1)^2+\cdots+(\mathrm{d}x^{n+1})^2\bigr)\\
&=(\mathrm{d}u^1)^2+\cdots+(\mathrm{d}u^n)^2+\mathrm{d}f^2
\end{align*}
给出。例如，$S^2$ 的上半球由映射 $X(u,v)=(u,v,\sqrt{1-u^2-v^2})$ 参数化。在这些坐标下，圆度量可写成
\begin{align*}
\bar g=X^*\bar g&=\mathrm{d}u^2+\mathrm{d}v^2+\left(\frac{u\,\mathrm{d}u+v\,\mathrm{d}v}{\sqrt{1-u^2-v^2}}\right)^2\\
&=\frac{(1-v^2)\,\mathrm{d}u^2+(1-u^2)\,\mathrm{d}v^2+2uv\,\mathrm{d}u\,\mathrm{d}v}{1-u^2-v^2}.
\end{align*}
\end{example}

\begin{example}[旋转曲面上的诱导度量]\label{example:ch13-revolution}
设 $C$ 是半平面 $\{(r,z):r>0\}$ 中的一个嵌入一维子流形，$S_C$ 是如例\ref{example:ch05-revolution}所述由 $C$ 生成的旋转曲面。为计算 $S_C$ 上的诱导度量，任取 $C$ 的一个光滑局部参数化 $\gamma(t)=(a(t),b(t))$，注意到映射 $X(t,\theta)=(a(t)\cos\theta,a(t)\sin\theta,b(t))$ 给出 $S_C$ 的一个光滑局部参数化（只要把 $(t,\theta)$ 限制在平面的一个充分小的开子集上）。于是我们可以算出
\begin{align*}
X^*\bar g&=\mathrm{d}(a(t)\cos\theta)^2+\mathrm{d}(a(t)\sin\theta)^2+\mathrm{d}(b(t))^2\\
&=(a'(t)\cos\theta\,\mathrm{d}t-a(t)\sin\theta\,\mathrm{d}\theta)^2\\
&\quad+(a'(t)\sin\theta\,\mathrm{d}t+a(t)\cos\theta\,\mathrm{d}\theta)^2+(b'(t)\,\mathrm{d}t)^2\\
&=(a'(t)^2+b'(t)^2)\,\mathrm{d}t^2+a(t)^2\,\mathrm{d}\theta^2.
\end{align*}
特别地，若 $\gamma$ 是一条单位速率曲线，即 $|\gamma'(t)|^2=a'(t)^2+b'(t)^2=1$，则它简化为简单公式 $\mathrm{d}t^2+a(t)^2\,\mathrm{d}\theta^2$。

下面是几个熟悉的旋转曲面例子。
\begin{enumerate}[(a)]
\item 例\ref{example:ch05-revolution}中描述的嵌入环面是由圆 $(r-2)^2+z^2=1$ 生成的旋转曲面。取该圆的单位速率参数化 $\gamma(t)=(2+\cos t,\sin t)$，便得到诱导度量的公式 $\mathrm{d}t^2+(2+\cos t)^2\,\mathrm{d}\theta^2$。
\item 单位球面（去掉南北两极）是一张旋转曲面，其生成曲线是由 $\gamma(t)=(\sin t,\cos t)$（$0<t<\pi$）参数化的半圆。诱导度量为 $\mathrm{d}t^2+\sin^2 t\,\mathrm{d}\theta^2$。
\item 单位圆柱面 $x^2+y^2=1$ 是一张旋转曲面，其生成曲线是由 $\gamma(t)=(1,t)$（$t\in\mathbb{R}$）参数化的竖直线。诱导度量为 $\mathrm{d}t^2+\mathrm{d}\theta^2$。
\end{enumerate}
\end{example}

再看上例最后一种情形。它表明，对圆柱面由 $X(t,\theta)=(\cos\theta,\sin\theta,t)$ 给出的每个局部参数化，诱导度量 $X^*\bar g$ 就是 $(t,\theta)$-平面上的 Euclid 度量。换句话说，对圆柱面上任意一点 $p$，适当限制 $X$ 便给出 $(\mathbb{R}^2,\bar g)$ 的某个开子集与圆柱面上 $p$ 的一个带诱导度量的邻域之间的一个 Riemann 等距。因此圆柱面上的诱导度量是平坦的。一个生活在圆柱面上的二维生物，单凭局部几何测量将无法把它周围的环境与 Euclid 平面区分开来。这说明一个度量是否平坦这个问题，有时会有出人意料的答案。

要发展出足够的工具来系统地判定哪些度量平坦、哪些不平坦，需要用到超出本书范围的技术。正如证明两个拓扑空间不同胚需要找到区分它们的拓扑不变量一样，为证明两个 Riemann 流形不局部等距，必须引入在 Riemann 等距下保持不变的局部不变量，并证明不同的度量具有不同的不变量。Riemann 度量的基本不变量称为它的\emph{曲率}；它是度量偏离平坦程度的一个定量刻画。关于 Riemann 曲率理论的介绍，可参见例如\cite{LeeRM}。

目前，我们只能满足于下一个命题，它用一种相当特设的方法回答了旋转曲面的这一问题。

\begin{proposition}[旋转曲面的平坦性判别法]\label{proposition:ch13-revolution-flat}
设 $C\subseteq H$ 是半平面 $H=\{(r,z):r>0\}$ 中一个连通的嵌入一维子流形，$S_C$ 是由 $C$ 生成的旋转曲面。则 $S_C$ 上的诱导度量是平坦的当且仅当 $C$ 是某条直线的一部分。
\end{proposition}

\begin{proof}
首先设 $C$ 是某条直线的一部分。则它有形如 $\gamma(t)=(Pt+K,Qt+L)$ 的参数化，其中 $P,Q,K,L$ 为常数，且 $P,Q$ 不全为零。通过重新标度 $t$ 变量，可设 $\gamma$ 是单位速率的。若 $Q=0$，则 $S_C$ 是平面 $z=L$ 的一个开子集，因而是平坦的。若 $P=0$，则 $S_C$ 是圆柱面 $x^2+y^2=K^2$ 的一部分，可以用与例\ref{example:ch13-revolution}(c)中单位圆柱面相同的方式证明它平坦。另一方面，若 $P,Q$ 都不为零，则 $S_C$ 是一张圆锥的一部分，例\ref{example:ch13-revolution}表明诱导度量为 $\mathrm{d}t^2+(Pt+K)^2\,\mathrm{d}\theta^2$。在任意一点的邻域内，坐标变换 $(u,v)=((t+K/P)\cos P\theta,(t+K/P)\sin P\theta)$ 把 Euclid 度量 $\mathrm{d}u^2+\mathrm{d}v^2$ 拉回到 $\mathrm{d}t^2+(Pt+K)^2\,\mathrm{d}\theta^2$，故此度量平坦。（设想把一张纸圆锥沿一侧剪开再摊平。）

反之，假设 $S_C$ 平坦，我们来证明 $C$ 是某条直线的一部分。设 $\gamma(t)=(a(t),b(t))$ 是 $C$ 的一个局部参数化。利用习题\ref{exercise:ch13-05}的结果，可设 $\gamma$ 是单位速率的，于是 $a'(t)^2+b'(t)^2=1$。如例\ref{example:ch13-revolution}，诱导度量为 $\mathrm{d}t^2+a^2\,\mathrm{d}\theta^2$。于是由
\begin{align*}
E^1=\frac{\partial}{\partial t},\qquad E^2=\frac1a\frac{\partial}{\partial\theta}
\end{align*}
给出的局部标架 $(E^1,E^2)$ 是标准正交的。任何其他标准正交标架 $(\widetilde E^1,\widetilde E^2)$ 都可写成
\begin{align*}
\widetilde E^1&=uE^1+vE^2=u\frac{\partial}{\partial t}+\frac va\frac{\partial}{\partial\theta},\\
\pm\widetilde E^2&=vE^1-uE^2=v\frac{\partial}{\partial t}-\frac ua\frac{\partial}{\partial\theta},
\end{align*}
其中 $u,v$ 是 $(t,\theta)$ 的光滑函数，满足 $u^2+v^2=1$。因为度量平坦，可以选取 $u,v$ 使 $(\widetilde E^1,\widetilde E^2)$ 成为一个交换标准正交标架（定理\ref{theorem:ch13-flatness}）。利用坐标下 Lie 括号的公式\eqref{align:ch08-lie-bracket-coord}，这意味着
\begin{align*}
0=[\widetilde E^1,\widetilde E^2]
&=\left(u\frac{\partial v}{\partial t}+\frac va\frac{\partial v}{\partial\theta}\right)\frac{\partial}{\partial t}
-\left(u\frac{\partial}{\partial t}\!\left(\frac ua\right)+\frac va\frac{\partial}{\partial\theta}\!\left(\frac ua\right)\right)\frac{\partial}{\partial\theta}\\
&\quad-\left(v\frac{\partial u}{\partial t}-\frac ua\frac{\partial u}{\partial\theta}\right)\frac{\partial}{\partial t}
-\left(v\frac{\partial}{\partial t}\!\left(\frac va\right)-\frac ua\frac{\partial}{\partial\theta}\!\left(\frac va\right)\right)\frac{\partial}{\partial\theta}.
\end{align*}
为简化此式，我们用简写记号 $f_\theta=\partial f/\partial\theta$、$f_t=\partial f/\partial t$。注意到 $u^2+v^2=1$ 蕴含 $uu_\theta+vv_\theta=uu_t+vv_t=0$，而 $a$ 只依赖于 $t$ 蕴含 $a_\theta=0$、$a_t=a'$。把这些关系代入上式并化简，得到
\begin{align*}
0=(uv_t-vu_t)\frac{\partial}{\partial t}+\left(\frac{a'-vu_\theta+uv_\theta}{a^2}\right)\frac{\partial}{\partial\theta},
\end{align*}
它蕴含
\begin{align}
uv_t-vu_t&=0,\label{align:ch13-13-2}\\
vu_\theta-uv_\theta&=a'.\label{align:ch13-13-3}
\end{align}
因为 $u^2+v^2\ne1$，每一点都有一个邻域使 $u$ 或 $v$ 非零。在任何 $v\ne0$ 的开子集上，\eqref{align:ch13-13-2} 表明 $u/v$ 关于 $t$ 的导数为零。于是可写成 $u=fv$，其中 $f$ 只是 $\theta$ 的函数。由 $u^2+v^2=1$ 得 $v^2(f^2+1)=1$，故 $v=\pm1/\sqrt{f^2+1}$ 也只是 $\theta$ 的函数，$u=\pm\sqrt{1-v^2}$ 亦然。在 $u\ne0$ 处有类似论证。但这样一来\eqref{align:ch13-13-3}表明 $a'$ 与 $t$ 无关，故为常数，从而 $b'=\pm\sqrt{1-(a')^2}$ 也是常数。于是 $a,b$ 都是 $t$ 的仿射函数，故 $C$ 的每一点都有一个邻域含于某条直线内。由于假设 $C$ 连通，便知整个 $C$ 含于同一条直线内。
\end{proof}

\begin{corollary}\label{corollary:ch13-sphere-not-flat}
$S^2$ 上的圆度量不是平坦的。
\end{corollary}

\subsection{法丛}

设 $(M,g)$ 是一个 $n$ 维带边或不带边的 Riemann 流形，$S\subseteq M$ 是一个 $k$ 维 Riemann 子流形（也带边或不带边）。正如我们对 $\mathbb{R}^n$ 的子流形所做的那样，对任意 $p\in S$，若一个向量 $v\in T_pM$ 关于内积 $\langle\cdot,\cdot\rangle_g$ 与 $T_pS$ 中每个向量都正交，就称它\emph{正交于} $S$。$S$ 在 $p$ 处的\emph{法空间}是 $T_pM$ 的子空间 $N_pS\subseteq T_pM$，由所有在 $p$ 处正交于 $S$ 的向量组成；$S$ 的\emph{法丛}是 $NS\subseteq TM$，由 $S$ 各点处所有法空间的并集组成。投影 $\pi_{NS}:NS\to S$ 定义为 $\pi:TM\to M$ 在 $NS$ 上的限制。下一个命题的证明与推论\ref{corollary:ch10-subbundle-criterion}相同。

\begin{proposition}[Riemann 子流形的法丛]\label{proposition:ch13-normal-bundle}
设 $(M,g)$ 是一个带边或不带边的 Riemann $n$ 流形。对任意带边或不带边的浸入 $k$ 维子流形 $S\subseteq M$，法丛 $NS$ 是 $TM|_S$ 的一个光滑秩 $(n-k)$ 子丛。对每个 $p\in S$，在 $p$ 的某个邻域上存在 $NS$ 的一个关于 $g$ 标准正交的光滑标架。
\end{proposition}

\begin{exercise}\label{exercise:ch13-22}
证明上一命题。
\end{exercise}
\end{document}
